\documentclass[12pt,a4paper]{article}
\usepackage[T2A]{fontenc}
\usepackage[utf8]{inputenc}
\usepackage[russian]{babel}
\usepackage{amsmath,amssymb,geometry,booktabs,array,enumitem,microtype}
\geometry{left=24mm,right=22mm,top=22mm,bottom=22mm}
\setlength{\parindent}{1.25cm}
\setlength{\parskip}{4pt}
\newcommand{\DTFT}{\mathcal{F}_{\mathrm{DT}}}
\newcommand{\dd}{\delta}
\begin{document}

\begin{center}
{\Huge\bfseries Тест № 1}\\[8pt]
{\large Основы цифровой обработки сигналов — подробные решения}
\end{center}
\vspace{8mm}
\section*{1. Спектр окна Ханна, начинающегося при $t=0$}
Дано окно Ханна длительностью $\tau$:
\[
w_H(t)=\begin{cases}
\dfrac12\left(1+\cos\dfrac{2\pi(t-\tau/2)}{\tau}\right),&0\leq t\leq\tau,\\[5pt]
0,&\text{иначе}.
\end{cases}
\]
Для симметричного относительно нуля окна той же длительности преобразование Фурье (лекция 1, слайд 15) равно
\[
W_0(f)=\frac{\sin(\pi f\tau)}{2\pi f(1-\tau^2 f^2)}.
\]
Заданное окно получается задержкой симметричного окна на $t_0=\tau/2$:
\[
w_H(t)=w_0(t-\tau/2).
\]
По теореме запаздывания
\[
w_0(t-t_0)\ \longleftrightarrow\ W_0(f)e^{-j2\pi f t_0}.
\]
Следовательно,
\[
\boxed{W_H(f)=\frac{\sin(\pi f\tau)}{2\pi f(1-\tau^2f^2)}e^{-j\pi f\tau}.}
\]
Значения в устранимых особых точках определяются по пределу.

\noindent\textbf{Правильный ответ: e.}

\section*{2. Спектр наблюдаемого участка гармонического сигнала}
Пусть
\[
x(t)=\cos(2\pi f_0t),\qquad f_0=200\ \text{Гц},\qquad \tau=0{,}1\ \text{с}.
\]
Наблюдение на интервале $0\leq t\leq\tau$ эквивалентно умножению на прямоугольное окно
\[
w(t)=\begin{cases}1,&0\leq t\leq\tau,\\0,&\text{иначе}.\end{cases}
\]
Его преобразование Фурье:
\begin{align*}
W(f)&=\int_0^\tau e^{-j2\pi ft}\,dt
=\tau\frac{\sin(\pi f\tau)}{\pi f\tau}e^{-j\pi f\tau}.
\end{align*}
Используем представление косинуса и свойство смещения спектра:
\[
\cos(2\pi f_0t)=\frac12e^{j2\pi f_0t}+\frac12e^{-j2\pi f_0t}.
\]
Таким образом, преобразование наблюдаемого сигнала $y(t)=x(t)w(t)$ равно
\[
\boxed{Y(f)=\frac12W(f-f_0)+\frac12W(f+f_0),}
\]
где
\[
W(f)=\tau\frac{\sin(\pi f\tau)}{\pi f\tau}e^{-j\pi f\tau}.
\]
Фазовый множитель обязателен, поскольку окно начинается при $t=0$, а не симметрично относительно нуля.

\noindent\textbf{Правильный ответ: b.}

\section*{3а. Спектр суммы двух гармоник}
Задан сигнал
\[
x(t)=3\cos(2\pi f_1t)+\cos(2\pi f_2t),\qquad
f_1=100\ \text{Гц},\quad f_2=200\ \text{Гц}.
\]
Для гармоники справедливо:
\[
A\cos(2\pi f_0t)\ \longleftrightarrow\
\frac A2\delta(f-f_0)+\frac A2\delta(f+f_0).
\]
По линейности преобразования Фурье
\[
\boxed{\begin{aligned}
X(f)={}&\frac32\delta(f-100)+\frac32\delta(f+100)\\
&+\frac12\delta(f-200)+\frac12\delta(f+200).
\end{aligned}}
\]
Следовательно, линии на частотах $\pm100$ Гц имеют веса $3/2$, а линии на $\pm200$ Гц --- веса $1/2$.

\noindent\textbf{Правильный ответ: c.}

\section*{3б. Спектр суммы гармоник после умножения на окно Ханна}
Пусть $y(t)=x(t)w_H(t)$, где $w_H(t)$ --- симметричное относительно нуля окно Ханна длительностью $\tau=0{,}1$ с.
Его спектр, как установлено в лекции 1, равен
\[
W_H(f)=\frac{\sin(\pi f\tau)}{2\pi f(1-\tau^2f^2)}.
\]
По теореме о спектре произведения
\[
Y(f)=X(f)*W_H(f).
\]
Используем фильтрующее свойство дельта-функции:
\[
\delta(f-f_0)*W_H(f)=W_H(f-f_0).
\]
Подставляя спектр $X(f)$ из пункта 3а, получаем
\[
\boxed{\begin{aligned}
Y(f)={}&\frac32W_H(f-f_1)+\frac32W_H(f+f_1)\\
&+\frac12W_H(f-f_2)+\frac12W_H(f+f_2),
\end{aligned}}
\]
где $f_1=100$ Гц, $f_2=200$ Гц.
Таким образом, каждая исходная спектральная линия заменяется сдвинутой копией спектра окна Ханна с соответствующим коэффициентом.

\noindent\textbf{Правильный ответ: b.}

\section*{Итог}
\[
\boxed{\text{1 --- e;\qquad 2 --- b;\qquad 3а --- c;\qquad 3б --- b.}}
\]

\clearpage
\begin{center}
{\Huge\bfseries Тест № 2}\\[8pt]
{\large Основы цифровой обработки сигналов — подробные решения}
\end{center}
\vspace{8mm}
\section*{Задача № 1. ДВПФ последовательности из 16 отсчётов}
\subsection*{Пункт а): спектр и его модуль}
Пусть $\dd[k]$ обозначает дискретный единичный импульс. Тогда
\[
 x[k]=\sum_{m=0}^{15}\dd[k-m]
 =\begin{cases}1,&0\le k\le15,\\0,&\text{иначе}.\end{cases}
\]
По определению ДВПФ в нормированной частоте $\nu$:
\[
 X(\nu)=\sum_{k=-\infty}^{\infty}x[k]e^{-j2\pi\nu k}
 =\sum_{k=0}^{15}e^{-j2\pi\nu k}
 =\frac{1-e^{-j32\pi\nu}}{1-e^{-j2\pi\nu}}.
\]
Используя $1-e^{-j2a}=2je^{-ja}\sin a$, получаем
\[
 \boxed{X(\nu)=e^{-j15\pi\nu}\frac{\sin(16\pi\nu)}{\sin(\pi\nu)}}.
\]
В точках целочисленных $\nu$ значения определяются предельным переходом.
Модуль спектра равен
\[
 \boxed{|X(\nu)|=\left|\frac{\sin(16\pi\nu)}{\sin(\pi\nu)}\right|}.
\]
Он периодичен с периодом 1, имеет максимум $|X(0)|=16$ и первые нули
при $\nu=\pm1/16$. Среди формул теста верен \textbf{вариант e}.
Графики \textbf{b и c} на предоставленном скриншоте оба удовлетворяют
этим критериям; однозначно отличить их по изображению нельзя.

\subsection*{Пункт б): модуляция косинусом}
Для $y_1[k]=x[k]\cos(2\pi k\,5/16)$ применим формулу Эйлера:
\[
 \cos(2\pi\nu_0k)=\frac12\bigl(e^{j2\pi\nu_0k}+e^{-j2\pi\nu_0k}\bigr),
 \quad\nu_0=\frac5{16}.
\]
По свойству сдвига ДВПФ $x[k]e^{j2\pi\nu_0k}\longleftrightarrow X(\nu-\nu_0)$:
\[
 \boxed{Y_1(\nu)=\frac12X\left(\nu-\frac5{16}\right)
 +\frac12X\left(\nu+\frac5{16}\right)}.
\]
\textbf{Ответ: c.}

\subsection*{Пункт в): модуляция синусом}
Для $y_2[k]=x[k]\sin(2\pi k\,5/16)$ используем
\[
 \sin(2\pi\nu_0k)=\frac{e^{j2\pi\nu_0k}-e^{-j2\pi\nu_0k}}{2j}.
\]
Следовательно,
\[
 \boxed{Y_2(\nu)=\frac1{2j}X\left(\nu-\frac5{16}\right)
 -\frac1{2j}X\left(\nu+\frac5{16}\right)}.
\]
\textbf{Ответ: c.}

\section*{Задача № 2. Окно Блэкмана}
\subsection*{Пункт а): ДВПФ окна}
При $N=16$ окно определено как
\[
 w_B[k]=\begin{cases}
 0.42-0.5\cos(2\pi k/16)+0.08\cos(4\pi k/16),&0\le k\le15,\\
 0,&\text{иначе}.
 \end{cases}
\]
Обозначим через $r[k]$ прямоугольное окно длиной 16 и через
$W(\nu)=e^{-j15\pi\nu}\sin(16\pi\nu)/\sin(\pi\nu)$ его ДВПФ.
Тогда
\[
 w_B[k]=r[k]\bigl(0.42-0.5\cos(2\pi k/16)+0.08\cos(4\pi k/16)\bigr).
\]
По линейности и теореме частотного сдвига:
\begin{align*}
 \boxed{W_B(\nu)={}}&0.42W(\nu)
 -0.25W\left(\nu-\frac1{16}\right)
 -0.25W\left(\nu+\frac1{16}\right)\\
 &+0.04W\left(\nu-\frac18\right)
 +0.04W\left(\nu+\frac18\right).
\end{align*}
\textbf{Ответ: c.} Коэффициенты $0.25$ и $0.04$ появляются
из разложения косинусов на пару комплексных экспонент.

\subsection*{Пункт б): окно Блэкмана, умноженное на косинус}
Для $y[k]=w_B[k]\cos(2\pi k\,5/16)$ снова применяем частотный сдвиг:
\[
 \boxed{Y(\nu)=\frac12W_B\left(\nu-\frac5{16}\right)
 +\frac12W_B\left(\nu+\frac5{16}\right)}.
\]
\textbf{Ответ: a.}

\section*{Задача № 3. ДВПФ дискретизированного гармонического сигнала}
Пусть $x(t)=\cos(2\pi f_0t)$, $f_0=17$~Гц, $f_s=170$~Гц.
Нормированная частота составляет
\[
 \nu_0=\frac{f_0}{f_s}=\frac{17}{170}=\frac1{10},\qquad
 x[k]=\cos\left(2\pi\frac{k}{10}\right).
\]
ДВПФ комплексной экспоненты содержит периодически повторяющиеся
дельта-функции. Разлагая косинус по Эйлеру, получаем
\[
 \boxed{X(\nu)=\frac12\sum_{n=-\infty}^{\infty}
 \left[\dd\left(\nu-\frac1{10}-n\right)
 +\dd\left(\nu+\frac1{10}-n\right)\right]}.
\]
Спектральные линии с весами $1/2$ лежат на $\nu=n\pm0.1$,
$n\in\mathbb Z$, то есть спектр периодичен с периодом 1.
\textbf{Ответы: график b; формула b.}

\section*{Задача № 4. Наложение спектров}
Заданы частоты $f_1=20$~кГц и $f_2=80$~кГц при дискретизации
с $f_s=60$~кГц. Частота Найквиста равна
\[
 f_N=\frac{f_s}{2}=30\text{ кГц}.
\]
Первая частота ($20$~кГц) лежит ниже частоты Найквиста,
поэтому алиасинг для первой гармоники отсутствует. Вторая ($80$~кГц)
лежит выше частоты Найквиста и отображается в полосу Найквиста
на частоту $|80-60|=20$~кГц.
Для косинусоид с нулевой фазой
\[
 \cos\left(2\pi\frac{80}{60}k\right)
 =\cos\left(2\pi\frac{20}{60}k\right),\qquad k\in\mathbb Z.
\]
Таким образом, по частоте две исходные аналоговые гармоники
невозможно однозначно различить по отсчётам без дополнительных
сведений. Для синусов с одинаковой известной фазой $0$ отсчёты
отличаются знаком; это фазовое различие не позволяет однозначно
определить исходную аналоговую частоту при неизвестной фазе.
\textbf{Ответ в тесте: a.}

\section*{Сводка ответов}
\begin{center}
\begin{tabular}{ll}
\toprule
Пункт & Ответ\\
\midrule
1а, график & b и c визуально соответствуют формуле\\
1а, формула & e\\
1б & c\\
1в & c\\
2а & c\\
2б & a\\
3, график & b\\
3, формула & b\\
4 & a\\
\bottomrule
\end{tabular}
\end{center}

\clearpage
\begin{center}
{\Huge\bfseries Тест № 3}\\[8pt]
{\large Основы цифровой обработки сигналов — подробные решения}
\end{center}
\vspace{8mm}
\section*{Задача № 1. Прямоугольное окно}
Дано окно длиной $N=8$:
\[
 w[k]=\begin{cases}1,&0\le k\le7,\\0,&\text{иначе}.\end{cases}
\]
Его дискретное во времени преобразование Фурье (ДВПФ):
\begin{align*}
W(\nu)&=\sum_{k=0}^{7}e^{-j2\pi\nu k}
=\frac{1-e^{-j16\pi\nu}}{1-e^{-j2\pi\nu}}\\
&=e^{-j7\pi\nu}\frac{\sin(8\pi\nu)}{\sin(\pi\nu)}.
\end{align*}
В точках устранимых особенностей значение определяется по пределу, в частности $W(0)=8$. Следовательно,
\[
|W(\nu)|=\left|\frac{\sin(8\pi\nu)}{\sin(\pi\nu)}\right|.
\]
\textbf{Ответ на вопрос о формуле: вариант b.}

Для восьмиточечного ДПФ $\nu_n=n/8$; $W_8[0]=8$, а $W_8[1],\ldots,W_8[7]=0$. После дополнения исходной последовательности восемью нулями шестнадцатиточечное ДПФ берёт отсчёты той же ДВПФ на более частой сетке $\nu_n=n/16$, поэтому видны ненулевые отсчёты боковых лепестков.

\textbf{Ответ на сопоставление: верхний график --- пункт б; нижний --- пункт а.}

\section*{Задача № 2. Шестнадцатиточечное ДПФ}
Используем ненормированное прямое ДПФ:
\[
X[n]=\sum_{k=0}^{15}x[k]e^{-j2\pi nk/16}.
\]
\subsection*{а) Шестнадцать единичных отсчётов}
Для $x[k]=1$, $k=0,\ldots,15$, по ортогональности экспонент
\[
X[n]=\begin{cases}16,&n=0,\\0,&n=1,\ldots,15.\end{cases}
\]
Это \textbf{первая таблица и третий набор графиков}.
\subsection*{б) Косинус}
\[
y_1[k]=\cos\left(2\pi\frac5{16}k\right)
=\frac12 e^{j2\pi5k/16}+\frac12e^{-j2\pi5k/16}.
\]
По ортогональности экспонент
\[
Y_1[n]=\begin{cases}8,&n\in\{5,11\},\\0,&\text{иначе}.\end{cases}
\]
Это \textbf{третья таблица и первый набор графиков}.
\subsection*{в) Синус}
\[
y_2[k]=\sin\left(2\pi\frac5{16}k\right)
=\frac1{2j}\left(e^{j2\pi5k/16}-e^{-j2\pi5k/16}\right).
\]
Отсюда
\[
Y_2[n]=\begin{cases}-8j,&n=5,\\8j,&n=11,\\0,&\text{иначе}.\end{cases}
\]
Это \textbf{вторая таблица и второй набор графиков}.

\section*{Задача № 3. Окно Блэкмана}
При $N=16$ и $0\le k\le15$:
\[
w_B[k]=0.42-0.5\cos\left(2\pi\frac{k}{N}\right)
+0.08\cos\left(2\pi\frac{2k}{N}\right).
\]
Разложение косинусов на комплексные экспоненты и ортогональность базиса ДПФ дают
\[
W_B[n]=\begin{cases}
0.42N,&n=0,\\
-0.25N,&n=1,15,\\
0.04N,&n=2,14,\\
0,&\text{иначе}.
\end{cases}
\]
\textbf{Ответ № 3а: вариант c.}

Для $y[k]=w_B[k]\cos(2\pi5k/16)$ используем свойство смещения спектра:
\[
Y[n]=\frac12W_B[(n-5)\bmod16]+\frac12W_B[(n+5)\bmod16].
\]
Подстановка ненулевых коэффициентов даёт
\[
Y[n]=\begin{cases}
0.21N,&n=5,11,\\
-0.125N,&n=4,6,10,12,\\
0.02N,&n=3,7,9,13,\\
0,&\text{иначе}.
\end{cases}
\]
Все коэффициенты действительны. \textbf{Ответ № 3б: график a.}

\section*{Задача № 4. Циклический сдвиг}
Исходная четырёхточечная последовательность: $x[k]=[1,2,3,4]$.
Дано
\[
Y[n]=e^{-j2\pi n/4}X[n].
\]
По теореме циклического запаздывания
\[
x[(k-m)\bmod N]\;\longleftrightarrow\;X[n]e^{-j2\pi nm/N}.
\]
При $m=1$, $N=4$:
\[
y[k]=x[(k-1)\bmod4]=[4,1,2,3].
\]
\textbf{Ответ № 4: график b.}

\section*{Итог ответов}
\begin{center}
\begin{tabular}{ll}
\toprule
Пункт & Ответ\\
\midrule
№ 1, графики & верхний --- б; нижний --- а\\
№ 1, формула ДВПФ & b\\
№ 2, таблицы сверху вниз & а, в, б\\
№ 2, графики сверху вниз & б, в, а\\
№ 3а & c\\
№ 3б & a\\
№ 4 & b\\
\bottomrule
\end{tabular}
\end{center}

\clearpage
\begin{center}
{\Huge\bfseries Тест № 4}\\[8pt]
{\large Основы цифровой обработки сигналов — подробные решения}
\end{center}
\vspace{8mm}
\section*{Вопрос 1. Минимальное затухание фильтра Баттерворта}
\textbf{Дано:} 16-разрядный биполярный АЦП, аналоговый ФНЧ Баттерворта порядка 3, $f_c=10$ кГц. Требуется определить минимальное затухание в полосе подавления в дБ, округлённое до целого.

В соответствии с моделью лекции шум квантования идеального $n$-разрядного АЦП при полномасштабной синусоиде оценивается отношением сигнал/шум
\[
\mathrm{SNR}_{\mathrm{dB}}=6.02n+1.76.
\]
Для $n=16$ имеем
\[
\mathrm{SNR}_{\mathrm{dB}}=6.02\cdot16+1.76=98.08\,\text{дБ}.
\]
В рамках принятой в лекции методики требуемое минимальное ослабление вне рабочей полосы выбирается равным этому значению:
\[
A_{\min}=98.08\,\text{дБ}.
\]
\textbf{Ответ: $\boxed{98\,\text{дБ}}$.}

\section*{Вопрос 2. Минимальная частота дискретизации}
\textbf{Дано:} те же параметры, что и в вопросе 1: $f_c=10$ кГц, порядок Баттерворта $n=3$, $A_{\min}=98.08$ дБ.

Квадрат модуля частотной характеристики фильтра:
\[
|H(f)|^2=\frac{1}{1+(f/f_c)^{2n}}=\frac{1}{1+(f/f_c)^6}.
\]
Ослабление в дБ определяется выражением
\[
A(f)=-10\log_{10}|H(f)|^2
=10\log_{10}\!\left(1+\left(\frac{f}{f_c}\right)^6\right).
\]
На границе полосы подавления $f_{\max}$ положим $A(f_{\max})=A_{\min}$. Тогда
\begin{align*}
10\log_{10}\!\left(1+\left(\frac{f_{\max}}{f_c}\right)^6\right)&=98.08,\\
\left(\frac{f_{\max}}{f_c}\right)^6&=10^{98.08/10}-1,\\
f_{\max}&=f_c\left(10^{98.08/10}-1\right)^{1/6}.
\end{align*}
Подстановка $f_c=10$ кГц даёт
\[
f_{\max}=10\left(10^{9.808}-1\right)^{1/6}\approx431.188\,\text{кГц}.
\]
По критерию, принятому в лекции, частота Найквиста должна быть не ниже этой границы:
\[
\frac{f_{s,\min}}2\ge f_{\max}.
\]
Следовательно,
\[
f_{s,\min}=2f_{\max}\approx862.376\,\text{кГц}.
\]
\textbf{Ответ: $\boxed{862\,\text{кГц}}$ (округление до целого).}

\section*{Вопрос 3. Ослабление фильтра Чебышёва I рода}
\textbf{Дано:} порядок $n=3$, $f_c=40$ Гц, пульсации в полосе пропускания $R_p=0.5$ дБ. Найти положительную величину затухания на $f=2f_c$.

Коэффициент передачи мощности:
\[
|H(\nu)|^2=\frac{1}{1+\varepsilon^2T_n^2(\nu)},\qquad \nu=\frac{f}{f_c}.
\]
Из определения уровня пульсаций:
\[
R_p=10\log_{10}(1+\varepsilon^2),\qquad
\varepsilon^2=10^{R_p/10}-1=10^{0.05}-1\approx0.12201845.
\]
Полином Чебышёва третьего порядка:
\[
T_3(\nu)=4\nu^3-3\nu,\qquad T_3(2)=4\cdot2^3-3\cdot2=26.
\]
Искомое положительное ослабление:
\begin{align*}
A(2f_c)&=-10\log_{10}|H(2)|^2\\
&=10\log_{10}\left(1+\varepsilon^2T_3^2(2)\right)\\
&=10\log_{10}\left(1+(10^{0.05}-1)\cdot26^2\right)\\
&\approx19.216\,\text{дБ}.
\end{align*}
\textbf{Ответ: $\boxed{19\,\text{дБ}}$.}

\section*{Вопрос 4. Число уровней квантования}
\textbf{Дано:} разрядность $n=16$ бит.

Число цифровых кодов идеального $n$-разрядного АЦП:
\[
L=2^n=2^{16}=65536.
\]
\textbf{Ответ: $\boxed{65536}$.}

\section*{Вопрос 5. Шаг дискретизации}
\textbf{Дано:} частота дискретизации $f_s=44100$ Гц. Найти $\Delta t$ в мкс.

Период дискретизации обратен частоте:
\[
\Delta t=\frac1{f_s}=\frac1{44100}\,\text{с}
=\frac{10^6}{44100}\,\text{мкс}
\approx22.6757\,\text{мкс}.
\]
\textbf{Ответ: $\boxed{23\,\text{мкс}}$.}

\section*{Вопрос 6. Длительность несжатой аудиозаписи}
\textbf{Дано:} моно, глубина кодирования 16 бит/отсчёт, $f_s=44100$ Гц, размер файла 280 КБ. Требуется оценка времени в целых секундах.

Один отсчёт одного канала занимает
\[
b=\frac{16}{8}=2\,\text{байта}.
\]
Скорость потока несжатых аудиоданных:
\[
R=f_s\,b\,C=44100\cdot2\cdot1=88200\,\text{байт/с},
\]
где $C=1$ --- число каналов. Если $1\,\text{КБ}=1024$ байта, оценка длительности без учёта небольшого заголовка WAV:
\[
T\approx\frac{280\cdot1024}{88200}=3.2508\,\text{с}.
\]
Если считать $1\,\text{КБ}=1000$ байт, получаем $T\approx3.1746$ с. В обоих случаях округление до целой секунды даёт одинаковый результат.

\textbf{Ответ: $\boxed{3\,\text{с}}$.}

\section*{Итоговая таблица}
\begin{center}
\begin{tabular}{cll}
\toprule
Вопрос & Величина & Ответ \\
\midrule
1 & Минимальное затухание & 98 дБ \\
2 & Минимальная частота дискретизации & 862 кГц \\
3 & Ослабление Чебышёва на $2f_c$ & 19 дБ \\
4 & Уровни квантования & 65536 \\
5 & Шаг дискретизации & 23 мкс \\
6 & Длительность записи & 3 с \\
\bottomrule
\end{tabular}
\end{center}
\noindent\textit{Примечание.} Вопросы 1--2 решены в соответствии с учебной методикой лекции: требуемое подавление фильтра отождествляется с идеальным SNR АЦП. Для проектирования реального антиалиасингового фильтра дополнительно учитываются уровни внеполосных сигналов, полоса полезного сигнала и допустимый уровень алиасинга.

\clearpage
\begin{center}
{\Huge\bfseries Тест № 5}\\[8pt]
{\large Основы цифровой обработки сигналов — подробные решения}
\end{center}
\vspace{8mm}
\section*{Условие}
Дан бесконечный гармонический сигнал
\[
x(t)=\sin(2\pi f_0t),\qquad \Delta t=\frac{1}{10f_0}.
\]
Отсчёты $x(k\Delta t)$ удерживаются постоянными на каждом интервале
$[k\Delta t,(k+1)\Delta t)$. Требуется найти спектр $Y(f)$ выходного
сигнала фиксатора и выбрать правильный ответ.

\section*{Решение}
\subsection*{1. Спектр исходного сигнала}
Из формулы Эйлера
\[
\sin(2\pi f_0t)=\frac{e^{j2\pi f_0t}-e^{-j2\pi f_0t}}{2j}
\]
следует
\[
X(f)=\frac{1}{2j}\bigl[\delta(f-f_0)-\delta(f+f_0)\bigr].
\]
Частота дискретизации $f_s=1/\Delta t=10f_0$.

\subsection*{2. Спектр последовательности отсчётов}
В соответствии с нормировкой, используемой в лекции,
\[
x_d(t)=\Delta t\sum_{k=-\infty}^{\infty}x(k\Delta t)\delta(t-k\Delta t).
\]
Её спектр представляет собой периодическое повторение исходного:
\[
X_d(f)=\sum_{n=-\infty}^{\infty}X(f-nf_s)
=\frac{1}{2j}\sum_{n=-\infty}^{\infty}
\left[\delta(f-f_0-nf_s)-\delta(f+f_0-nf_s)\right].
\]

\subsection*{3. Частотная характеристика фиксатора}
Для удержания значения в течение $\Delta t$ импульсная характеристика
равна $h(t)=1$ при $0\le t<\Delta t$ и нулю иначе. Поэтому
\[
H(f)=\int_0^{\Delta t}e^{-j2\pi ft}\,dt
=\frac{\sin(\pi f\Delta t)}{\pi f}e^{-j\pi f\Delta t}.
\]
Формально умножение $X_d(f)H(f)$ при указанной нормировке $x_d$
даёт дополнительный множитель $\Delta t$ у физически ожидаемого
выходного спектра. Чтобы описать именно выход удерживающего устройства
$y(t)=x(k\Delta t)$ на каждом интервале, следует использовать
ненормированную импульсную последовательность
$\sum_k x(k\Delta t)\delta(t-k\Delta t)$; её спектр равен
$f_s X_d(f)$. Отсюда
\[
Y(f)=f_sX_d(f)H(f)=X_d(f)G(f),\quad
G(f)=\frac{\sin(\pi f/f_s)}{\pi f/f_s}e^{-j\pi f/f_s}.
\]

\subsection*{4. Итоговое выражение}
Положим $a_n=f_0+nf_s=f_0(1+10n)$. Обозначим
\[
G(f)=e^{-j\pi f/f_s}\frac{\sin(\pi f/f_s)}{\pi f/f_s}.
\]
Тогда точное выражение для спектра выходного сигнала:
\[
\boxed{\displaystyle
Y(f)=\frac{1}{2j}\sum_{n\in\mathbb Z}
\left[G(f_0+nf_s)\delta(f-f_0-nf_s)
-G(-f_0+nf_s)\delta(f+f_0-nf_s)\right].}
\]
Для первой группы линий коэффициент равен
\[
G(a_n)=e^{-j\pi(0.1+n)}
\frac{\sin(\pi(0.1+n))}{\pi(0.1+n)}.
\]
Для второй группы нужно подставить в $G$ отрицательную
частоту $-f_0+nf_s$. Её коэффициенты комплексно сопряжены
соответствующим коэффициентам первой ветви. Именно такое
представление гарантирует действительность $y(t)$.

\section*{Ответ}
\textbf{Вариант а.} Для исходной синусоиды спектр содержит
разность двух групп спектральных составляющих; вариант б с плюсом
этому не соответствует. В формуле варианта а на изображении
коэффициент содержит $1/f_0$, тогда как для физически нормированного
выхода удерживающего ЦАП он должен быть безразмерным:
$\sin(\pi(0.1+n))/[\pi(0.1+n)]$.

\end{document}
